\documentclass[12pt]{article}

\usepackage[margin=1in]{geometry}
\usepackage{enumitem}
\usepackage{hyperref}

\title{Financial Engineering (F414) \\ Project Assignment}
\author{Universidad de San Andrés}
\date{2026}

\begin{document}

\maketitle

\section*{Overview}

\noindent The objective of this project is to design, implement, and evaluate a
\textbf{production-level quantitative trading strategy} based on one of the
strategy families studied during the course.

\noindent Students will work in groups to translate academic research into a realistic,
robust, and fully reproducible trading system.

\section*{Group Size}

\noindent Students must work in groups of \textbf{up to 4 members}.

\section*{Strategy Families}

\noindent Each group must select \textbf{one} of the following strategy families studied
during the course:

\begin{itemize}
\item Trend Following (Time-Series Momentum)
\item Momentum (Cross-Sectional Momentum)
\item Value
\item Quality
\item Low Volatility
\item Carry Trade
\item Volatility Risk Premium
\item Statistical Arbitrage
\end{itemize}

\section*{Literature Basis}

\noindent Your strategy must be grounded in academic research. Each group must rely on:

\begin{itemize}
\item The \textbf{main paper} discussed in class
\item The \textbf{complementary paper} assigned as homework
\item \textbf{Two additional academic papers} of your choice related to the
same strategy family
\end{itemize}

\noindent The goal is not merely to replicate the paper, but to build a realistic
trading strategy inspired by the literature.

\section*{Project Objectives}

\noindent Each group must:

\begin{enumerate}[label=\arabic*.]

\item \textbf{Design a Production-Level Strategy}

Construct a systematic trading strategy based on the chosen theme. This
includes:

\begin{itemize}
\item Economic intuition behind the anomaly
\item Quantitative signal construction
\item Portfolio construction rules
\item Risk management methodology
\end{itemize}

\item \textbf{Develop a Realistic Backtesting Framework}

Your backtesting environment must be designed to resemble a real trading
system. The framework must include:

\begin{itemize}
\item Data ingestion and preprocessing
\item Realistic transaction costs
\item Portfolio rebalancing rules
\item Performance metrics (Sharpe ratio, drawdown, turnover, etc.)
\end{itemize}

\item \textbf{Perform Robustness Testing}

Beyond the historical backtest, you must evaluate the robustness of the
strategy through additional tests, such as:

\begin{itemize}
\item Scenario simulations
\item Stress testing
\item Parameter sensitivity analysis
\item Out-of-sample testing
\end{itemize}

\item \textbf{Build a Usable Trading Framework}

Your system must be runnable by any user. The framework should include
either:

\begin{itemize}
\item A \textbf{Graphical User Interface (GUI)}, or
\item A \textbf{Command Line Interface (CLI)}
\end{itemize}

The system must allow users to:

\begin{itemize}
\item Run backtests
\item Monitor strategy performance
\item Connect to a broker account (e.g., IBKR, Robinhood, etc.)
\item Execute trades automatically
\item Log and track executed trades
\item Monitor the portfolio in real time
\end{itemize}

\end{enumerate}

\section*{Use of AI Tools}

\noindent Students are allowed to use AI-assisted development tools, including but not
limited to:

\begin{itemize}
\item ChatGPT
\item Claude
\item Claude Code
\item Cursor
\item Codex
\item OpenClaw
\end{itemize}

\noindent However, students must provide \textbf{transparency regarding their usage}.

You must submit either:

\begin{itemize}
\item The full conversation logs with the LLMs used, or
\item Screenshots documenting the interactions with the AI tools
\end{itemize}

\section*{Code Submission}

\noindent All code and materials must be uploaded to a dedicated branch in the
following GitHub repository:

\begin{center}
\url{https://github.com/fedeglan/UdeSA-Financial-Engineering-Course}
\end{center}

\noindent Each group will be assigned a branch corresponding to its
\textbf{group number}. All project materials must be submitted there,
including:

\begin{itemize}
\item Source code
\item Documentation
\item Data processing scripts
\item Backtesting framework
\item Instructions to run the system
\end{itemize}

\section*{Deadlines and Presentations}

\textbf{Final submission deadline:} \textbf{June 1st, 2026}.

\noindent All groups must have their complete project uploaded to the GitHub repository
by this date.

\noindent Following submission, each group will present their project during the
scheduled presentation sessions.

\noindent Each group will have \textbf{40 minutes} to:

\begin{itemize}
\item Present the final product
\item Demonstrate the system in action
\item Explain the strategy and its economic intuition
\item Describe the development process and design decisions
\item Discuss the robustness tests and performance results
\end{itemize}

\noindent There will be \textbf{four presentation dates}, with \textbf{two groups presenting
per session}:

\begin{itemize}
\item June 3rd, 2026
\item June 10th, 2026
\item June 17th, 2026
\item June 24th, 2026
\end{itemize}

\section*{Expected Deliverables}

\begin{itemize}
\item A working trading framework (GUI or CLI)
\item A reproducible backtesting pipeline
\item Documentation explaining the strategy and system architecture
\item Academic references used
\item Evidence of AI-assisted development (logs or screenshots)
\end{itemize}

\end{document}